\section{Discussion and Limitations}
\label{sec:discussion}

\subsection{What the Agentic-Time Claim Means}

The central claim is architectural: causal history, synchronization, and
conflict isolation can share one filesystem representation, so independent
agents need not coordinate every update through an explicit repository commit.
The artifact implements concurrent frontiers, disjoint-fragment merge, retained
strains, and hosted and peer transports.  The controlled multi-replica
scenarios in \cref{subsec:eval-mechanisms} exercise those mechanisms, including
third-replica discovery and resolution propagation, but do not establish how a
real fleet of agents behaves.  The evaluation contains no repository trace,
non-loopback LAN or WAN deployment, conflict-frequency study, or scaling curve
over agents, replicas, files, and retained versions.  We therefore claim a substrate for that
workload, not measured superiority to Git, worktrees, or existing synchronizers
for an end-to-end development task.  Such a study is the most important
remaining evaluation rather than a missing storage feature in RC1.

For the intended workload, the first missing network experiment is a LAN
experiment rather than necessarily a long-haul WAN experiment: colocated
agents and developer machines make 1- and 10-Gbit/s Ethernet the primary
operating range.  Remote and cloud-agent cases add links from roughly
10-Mbit/s constrained upstreams through 10-Gbit/s data-center connectivity.
A useful evaluation should sweep link rate, RTT, and loss independently while
varying writers and changed fragments, then report synchronization completion
time, transferred bytes, CPU cost, and conflict visibility latency.  These are
future evaluation parameters; the loopback TLS scenarios in this paper make no
claim about any point in that matrix.

\subsection{Security Boundary}

The strongest deployment claim is intentionally conditional.  End-to-end
confidentiality and authenticated content require GCM content, the always-GCM
metadata path, and WriterSet enforcement.  XTS content is confidential but
malleable; plaintext content has no cryptographic protection.  Neither mode may
inherit the full secure-SaaS claim merely because it uses the same transport.
Even in the secure profile, the service learns the account, UUIDs, fragment
indices and causal versions, version vectors, sizes, timing, and access
patterns.  \sfs{} is consequently client-side encrypted, not metadata-oblivious
or zero knowledge.

The hosted service is assumed honest-but-curious.  It may reject an
unauthorized record frontier, and GCM rejects modified ciphertext, but no
protocol in v12 detects a service that withholds a recent frontier, presents an
old but valid view, or forks two clients.  The signed writer set provides
authorization, not SUNDR-style fork consistency.  A transparency log,
client-to-client witness exchange, or signed monotone root chain could strengthen
that model, at additional state and availability cost.  Endpoint compromise,
root-key theft, traffic analysis, and denial of service are also outside the
current protection boundary.  No independent cryptographic audit has been
performed.

Key rotation is a forward boundary, not retroactive erasure.  A removed writer
may retain ciphertext and keys obtained before removal.  The live state and
unresolved strains are re-encrypted into the new epoch, but pre-rotation parent
history and pins are not carried across it.  Applications that require one
continuous, decryptable history across revocation need a different archival-key
policy.

\subsection{Granularity, Retention, and Durability}

Fixed positions make offset lookup constant time and let disjoint fragments
merge without a content-specific data type.  They deliberately give up the
insertion stability and cross-file deduplication of content-defined chunking.
Two concurrent edits to different bytes in one fragment still conflict, and a
size-band transition rewrites the complete logical unit under fresh dots and
loses delta reuse against its old geometry.  The artifact's 32-MiB append
regression records the bounded forward-growth case in
\cref{subsec:eval-mechanisms}; it does not measure latency or repeated growth,
shrink, or oscillation around a boundary.  Manifest construction and comparison
are linear in total unit and fragment metadata; only position lookup is
\(O(1)\).  Workloads dominated by insertion shifts or cross-file duplicates
may prefer content-defined chunks or an application CRDT.

History is bounded by policy and key epoch.  Named cuts pin required fragments
against ordinary retention, but they are not an eternal-retention promise.
Reclamation, retention expiry, defragmentation, and trim are explicit; the
artifact has no background cleaner whose aged steady state is evaluated here.
The three-region layout therefore needs long-duration fragmentation and
space-amplification measurements under repeated overwrite, snapshot, sync, and
rotation workloads.

The alternating authenticated header is an atomic publication point, not a
claim that every path is log-free or formally crash-safe.  The asynchronous
mount path uses a WAL and bounded in-place changes use undo records.  The
artifact contains crash-seam tests, but this paper reports no real power-cut or
systematic block-reordering campaign.  Offline scanning is a recovery aid; its
highest-address choice among unrelated orphan records is explicitly heuristic.

\subsection{Surfaces and Format Maturity}

The raw-device driver and FUSE adapter expose a broad useful VFS subset rather
than proven POSIX conformance.  Alias and hard-link accounting remain
incomplete, and the FUSE view reports simplified link counts.  A FUSE mount and
kernel mount must not open the same image concurrently because their lock
domains do not meet.  The Linux kernel path is the evaluated high-performance
encrypted path.  The FUSE encrypted-read deficit in
\cref{subsec:eval-fuse} is characterized but unresolved; it matters when the
portable mount is expected to carry high-throughput encrypted reads.
The profile decomposes \sfs{} itself but does not comparatively instrument
gocryptfs, so it cannot attribute the threefold gap to splice, copying, or one
particular scheduling decision.  Those remain hypotheses for a dedicated
buffer-ownership and pipeline experiment, not findings of this paper.

The wasm adapter currently demonstrates that the core engine can operate on an
in-memory image through a browser-shaped API.  RC1 does not include the
browser-target build and JavaScript end-to-end evidence or a persistent
OPFS/IndexedDB backend needed to call it a deployed web filesystem.  Likewise,
macFUSE and WinFsp integration points are not substitutes for measurements on
macOS and Windows.  These are important execution surfaces, but the paper keeps
them experimental.

Finally, the reader accepts exactly v12 and provides no migration path from
development formats.  A release can reasonably make this clean cut, but a
long-lived application VFS needs a documented compatibility policy, migrations,
and rollback tooling.  The published performance report preserves aggregate
cells but not raw samples and per-run health logs; a durable research artifact
should archive both and bind them to full source and environment manifests.
